Define qualitative robustness.

Robustness depends on the metric that is used on the space of probability measures.
Risk measures are not either robust or not robust, but more or less robust.
There exists a tradeoff between tail sensitivity and robustness.

Let \(\Omega = \mathbb{N}\), \(X_i(\omega) = \omega(i)\), and \(\mathcal{F} = \sigma(X_i : i \in \mathbb{N})\).
For any Borel probability measure \(\mu\) on \(\mathbb{R}\) denote
\[P_{\mu} := \mu^{\otimes \mathbb{N}}.\]
\(\mathcal{A} \subseteq \mathcal{M}_1\), the probability measures on \(\mathbb{R}\), with metric \(d_A\)
\(d_B\) metric on \(\mathcal{M}_1\)
\(\mathcal{R}_{\rho}\) is robust on \(\mathcal{A}\) with respect to \(d_A\), \(d_B\), if
\[\forall \mu \in \mathcal{A}, \epsilon > 0 \quad \exists \delta > 0, n_0 \in \mathbb{N} \quad \forall n \geq n_0: \nu \in \mathcal{A}, d_A(\mu, \nu) \leq \delta \Rightarrow d_B(P_{\mu} \circ \hat{P}_n^{-1}, P_{\nu} \circ \hat{P}_n^{-1}) \leq \epsilon\]
where \(\hat{P}_n := \mathcal{R}_{\rho}(\frac{1}{n} \sum_{k=1}^n \delta_{X_k})\).

How can the weak topology and the \( \varphi \)-weak topologies be metrized
and what is \( \varphi \)-Robustness?

A set \(\mathcal{A} \subset \mathcal{M}_1\) is called uniformly \(\psi\)-integrating, if
\[\lim_{M \rightarrow \infty} \sup_{\nu \in \mathcal{A}} \int_{\{\psi \geq M\}} \psi d\nu = 0.\]

The Prohorov \(\psi\)-metric is given by:
\[d_{\psi}(\mu, \nu) = d_{Proh}(\mu, \nu) + \left| \int \psi d\mu - \int \psi d\nu \right|, \quad \mu, \nu \in \mathcal{M}_1,\]
where the Prohorov metric is defined as
\[d_{Proh}(\mu, \nu) = \inf \{ \epsilon > 0 : \mu(A) \leq \nu(A^{\epsilon}) + \epsilon \text{ for all } A \in \mathcal{B}(\mathbb{R}) \}\]
with \(A^{\epsilon} = \{ x \in \mathbb{R} : \inf_{a \in A} |x-a| < \epsilon \}, \epsilon > 0\).

While the weak topology which is induced by the Prohorov metric is insensitive to the tails, for unbounded \(\psi\) the \(\psi\)-weak 
topology is finer and does not neglect deviations of measures in the tails.

Definition (\(\psi\)-Robustness):
A risk functional \(\mathcal{R}_{\rho}\) is called \(\psi\)-robust on \(\mathcal{M} \subset \mathcal{M}_1\), 
if it is robust with respect to \(d_{\psi}\) and \(d_{Proh}\) on every uniformly \(\psi\)-integrating set \(\mathcal{A} \subset \mathcal{M}\).

State equivalent characterizations of \( \varphi \)-robustness.

Let \(\rho: L^{\infty} \rightarrow \mathbb{R}\) be a distribution-based risk measure and 
\(\bar{\rho}: L^1 \rightarrow \mathbb{R} \cup \{+\infty\}\) its unique extension. 
Suppose that \(\Psi\) is a finite Young function that satisfies the \(\Delta_2\)-condition, and set 
\(\psi(x) = 1 + \Psi(|x|), x \in \mathbb{R}\).

Then the following conditions are equivalent:

a) \(\bar{\rho}\) is finite on the Orlicz heart \(H^{\Psi} := H^{\Psi}(\Omega, \mathcal{F}, P) = \{ X \in L^0 : E[\Psi(c|X|)] < \infty\) 
    for all \(c > 0 \}\).
b) The mapping \(\mathcal{R}_{\rho}: \mathcal{M}_{1,c} \rightarrow \mathbb{R}\) is continuous with respect to the \(\psi\)-weak topology.

If any of the conditions above is satisfied, the mapping 
\(\mathcal{R}_{\rho}: \mathcal{M}_{\psi}^1 = \mathcal{M}(H^{\Psi}) \rightarrow \mathbb{R}\) 
is well defined and continuous with respect to the \(\psi\)-weak topology.

Any of the conditions a) - b) is equivalent to the \(\psi\)-robustness of \(\mathcal{R}_{\rho}: \mathcal{M}_{1,c} \rightarrow \mathbb{R}\).

If any of the equivalent conditions is satisfied, the mapping \(\mathcal{R}_{\rho}: \mathcal{M}(H^{\Psi}) \rightarrow \mathbb{R}\) is
- well defined,
- continuous with respect to the \(\psi\)-weak topology, and
- \(\psi\)-robust.


How is the index of qualitative robustness defined?

The question goes on with, What is the largest possible
index for convex distribution-based risk measures? Can you specify examples of
risk measures that attain this smallest upper bound for convex risk measures?

\(L^p\)-spaces are particularly important. 
Letting \(\Psi_p(x) = |x|^p/p\) for \(0 < p < \infty\), \(x \in \mathbb{R}\), the \(\Delta_2\)-condition is always satisfied. 
This is a Young function if \(p \geq 1\). This is required for an application of Theorem 4.3.8.

Definition (Index of Qualitative Robustness)

Assume that \(\rho: L^{\infty} \rightarrow \mathbb{R}\) is a distribution-based risk measure. 
The associated index of qualitative robustness is given by
\[iqr(\rho) = \frac{1}{\inf \{ p \in (0, \infty) : \mathcal{R}_{\rho} \text{ is } \psi_p\text{-robust on } \mathcal{M}_{1,c} \}} \in [0, \infty]\]

Convex risk measures always have an index of qualitative robustness that is at most 1. 
In view of Theorem 4.3.8 & Theorem 4.3.9, the index of qualitative robustness can then be written as
\[iqr(\rho) = \frac{1}{\inf \{ p \in [1, \infty) : \bar{\rho} \text{ is finite on } L^p \}} \in [0, 1]\].

Average Value at Risk at level \( α ∈ (0, 1) \), AV@Rα , is a special case of a distortion risk measure \( \rho_g \) with
\( g(t) = (\frac{t}{\alpha} \wedge 1 \)
Since the right hand side derivative of \( g \) is bounded, one obtains that \( \text{iqr}(\text{AV@R}_{\alpha} = 1\).
